\documentclass[
	% -- opções da classe memoir --
	12pt,				% tamanho da fonte
	openright,			% capítulos começam em pág ímpar (insere página vazia caso preciso)
	oneside,			% para impressão em verso e anverso. Oposto a oneside
	a4paper,			% tamanho do papel. 
	english,			% idioma adicional para hifenização
	french,				% idioma adicional para hifenização
	spanish,			% idioma adicional para hifenização
	brazil				% o último idioma é o principal do documento
	]{abntex2}

% ---
% Pacotes básicos 
% ---
\usepackage{lmodern}			% Usa a fonte Latin Modern			
\usepackage[T1]{fontenc}		% Selecao de codigos de fonte.
\usepackage[utf8]{inputenc}		% Codificacao do documento (conversão automática dos acentos)
\usepackage{lastpage}			% Usado pela Ficha catalográfica
\usepackage{indentfirst}		% Indenta o primeiro parágrafo de cada seção.
\usepackage{color}				% Controle das cores
\usepackage{graphicx}			% Inclusão de gráficos
\usepackage{microtype} 			% para melhorias de justificação
\usepackage{multirow}
\usepackage{url}
\usepackage[output=pdf]{lib/logoLuizRosa}
\usepackage{supertabular}
\renewcommand{\UrlBreaks}{\do\/\do\a\do\b\do\c\do\d\do\e\do\f\do\g\do\h\do\i\do\j\do\k\do\l\do\m\do\n\do\o\do\p\do\q\do\r\do\s\do\t\do\u\do\v\do\w\do\x\do\y\do\z\do\A\do\B\do\C\do\D\do\E\do\F\do\G\do\H\do\I\do\J\do\K\do\L\do\M\do\N\do\O\do\P\do\Q\do\R\do\S\do\T\do\U\do\V\do\W\do\X\do\Y\do\Z}

% ---
		
% ---
% Pacotes adicionais, usados apenas no âmbito do Modelo Canônico do abnteX2
% ---
\usepackage{lipsum}				% para geração de dummy text
% ---

% ---
% Pacotes de citações
% ---
\usepackage[brazilian,hyperpageref]{backref}	 % Paginas com as citações na bibl
\usepackage[alf]{abntex2cite}	% Citações padrão ABNT

% ---
% Informações de dados para CAPA e FOLHA DE ROSTO
% ---
\titulo{Diretrizes e Normas para Trabalhos Acadêmicos}
\autor{Equipe da Escola Professor Luiz Rosa}
\local{Jundiaí}
\data{2014}
\orientador{Professores da Escola Prof. Luiz Rosa}
\coorientador{}
\instituicao{%
  Escola Professor Luiz Rosa}
\tipotrabalho{}
% O preambulo deve conter o tipo do trabalho, o objetivo, 
% o nome da instituição e a área de concentração 
\preambulo{Orientações para a apresentação escrita dos Projetos Integradores originalmente preparadas para a Escola Prof. Luiz Rosa.}
% ---


% ---
% Configurações de aparência do PDF final

% alterando o aspecto da cor azul
\definecolor{blue}{RGB}{41,5,195}

% informações do PDF
\makeatletter
\hypersetup{
     	%pagebackref=true,
		pdftitle={\@title}, 
		pdfauthor={\@author},
    	pdfsubject={\imprimirpreambulo},
	    pdfcreator={LaTeX with abnTeX2},
		pdfkeywords={abnt}{latex}{abntex}{abntex2}{trabalho acadêmico}, 
		colorlinks=false,       		% false: boxed links; true: colored links
    	linkcolor=blue,          	% color of internal 	links
    	citecolor=blue,        		% color of links to bibliography
    	filecolor=magenta,      		% color of file links
		urlcolor=blue,
		bookmarksdepth=4
}
\makeatother
% --- 

% --- 
% Espaçamentos entre linhas e parágrafos 
% --- 

% O tamanho do parágrafo é dado por:
\setlength{\parindent}{1.3cm}

% Controle do espaçamento entre um parágrafo e outro:
\setlength{\parskip}{0.2cm}  % tente também \onelineskip

% ---
% compila o indice
% ---
\makeindex
% ---

% ----
% Início do documento
% ----
\begin{document}

% Retira espaço extra obsoleto entre as frases.
\frenchspacing 

% ----------------------------------------------------------
% ELEMENTOS PRÉ-TEXTUAIS
% ----------------------------------------------------------
% \pretextual

% ---
% Capa
% ---
\imprimircapa
% ---

% ---
% Folha de rosto
% (o * indica que haverá a ficha bibliográfica)
% ---
\imprimirfolhaderosto*
% ---

% ---
% Inserir folha de aprovação
% ---

% Isto é um exemplo de Folha de aprovação, elemento obrigatório da NBR
% 14724/2011 (seção 4.2.1.3). Você pode utilizar este modelo até a aprovação
% do trabalho. Após isso, substitua todo o conteúdo deste arquivo por uma
% imagem da página assinada pela banca com o comando abaixo:
%
% \includepdf{folhadeaprovacao_final.pdf}
%

% ---
% Dedicatória
% ---
\begin{dedicatoria}
   \vspace*{\fill}
   \centering
   \noindent
   \textit{ Este trabalho é dedicado às crianças adultas que,\\
   quando pequenas, sonharam em se tornar cientistas.} \vspace*{\fill}
\end{dedicatoria}
% ---

% ---
% Agradecimentos
% ---
\begin{agradecimentos}
Agradecimentos à equipe da Escola Professor Luiz Rosa para a criação deste trabalho. Em especial, à Professora Daiane, que produziu a primeira versão deste documento. À Professora Joyce, que contribuiu sobremaneira com o entendimento da norma ABNT. 
Agradecimento à coordenação (Margarete, Valquíria, Viviane e Frederico). O apoio de vocês à proposta pedagógica permite o andamento deste trabalho.

\end{agradecimentos}
% ---

% ---
% Epígrafe
% ---
\begin{epigrafe}
    \vspace*{\fill}
	\begin{flushright}
		\textit{``Diante da vastidão do tempo e da imensidão do universo, é um imenso prazer para mim dividir um planeta e uma época com você."\\
		(Carl Sagan)}
	\end{flushright}
\end{epigrafe}
% ---

% ---
% RESUMOS
% ---

% resumo em português
\setlength{\absparsep}{18pt} % ajusta o espaçamento dos parágrafos do resumo
\begin{resumo}
%Escrever aqui o resumo do trabalho.
Este documento explicita as normas que devem ser obedecidas para a escrita dos trabalhos finais do Projeto Integrador. Todas as orientações aqui apresentadas são baseadas nas normas da Associação Brasileira de Normas Técnicas (ABNT) %\sigla{ABNT}{Associação Brasileira de Normas Técnicas} para escrita de trabalhos científicos.

 \textbf{Palavras-chaves}: projeto integrador, normas ABNT.
\end{resumo}

% resumo em inglês
\begin{resumo}[Abstract]
 \begin{otherlanguage*}{english}
   This is the english abstract.

   \vspace{\onelineskip}
 
   \noindent 
   \textbf{Key-words}: latex. abntex. text editoration.
 \end{otherlanguage*}
\end{resumo}



% ---

% ---
% inserir lista de ilustrações
% ---
\pdfbookmark[0]{\listfigurename}{lof}
\listoffigures*
\cleardoublepage
% ---

% ---
% inserir lista de tabelas
% ---
\pdfbookmark[0]{\listtablename}{lot}
\listoftables*
\cleardoublepage
% ---

%\listadesiglas

% ---
% inserir o sumario
% ---
\pdfbookmark[0]{\contentsname}{toc}
\tableofcontents*
\cleardoublepage
% ---

% ----------------------------------------------------------
% ELEMENTOS TEXTUAIS
% ----------------------------------------------------------
\textual

% ----------------------------------------------------------
% Introdução (exemplo de capítulo sem numeração, mas presente no Sumário)
% ----------------------------------------------------------
\chapter*[Introdução]{Introdução}
\addcontentsline{toc}{chapter}{Introdução}
% ----------------------------------------------------------

\par Este roteiro pretende ser um instrumento facilitador da redação de um trabalho acadêmico / científico, o qual, por sua vez, apresenta rigor na sua forma organizacional bem como na sua apresentação racional e formal, seguindo as normas da ABNT.
\par É importante o esclarecimento sobre o objetivo, da Escola Professor Luiz Rosa que não é explicar como é feita a pesquisa científica em suas várias modalidades, mas sim apresentar orientações metodológicas para a elaboração de trabalhos acadêmicos denominados Projetos Integradores, originalmente concebidos para a Escola Prof. Luiz Rosa em seus cursos médio e técnico.
\par Um trabalho científico tem por objetivo divulgar informações desejadas e servir como referência a quem se interessa pelo assunto. Assim, a comunicação escrita é a forma mais adequada e privilegiada na difusão do conhecimento pesquisado e cientificamente produzido.
\par Essa produção escrita deve apresentar linguagem objetiva, descritiva, coerente e coesa no intuito de se tornar uma referência de pesquisa a outros estudiosos do assunto. É por isso que se dá importância à redação e à linguagem crítica e analítica do trabalho.
\par No \autoref{ch:formato}, contempla-se a formatação do trabalho acadêmico. A descrição dos elementos que compõem o Projeto Integrador está delineada no \autoref{ch:estrutura}.  Já no \autoref{ch:elementos}, há uma explicação sobre cada um desses elementos em um trabalho acadêmico, com especial ênfase aos elementos pré-textuais e suas características.
\par A preocupação com a produção escrita propriamente dita, relacionada à pesquisa, à análise de dados e às considerações finais, denominada elementos textuais, encontra-se discutida no \autoref{ch:pesquisa}, espaço em que se procura enfatizar a necessidade da escrita própria desvinculada de cópia pura de textos de outros autores, o que caracteriza a apropriação indevida de material alheio. Aprecia-se aqui a utilização de mecanismos norteadores da referência a outros autores.
\par O \autoref{ch:posTextuais} traz os elementos pós-textuais, que são a indicação dos materiais pesquisados, cujo objetivo é a documentação, completa ou parcial, para melhor esclarecimento do texto do Projeto Integrador.
\par Espera-se que estas orientações sejam úteis àqueles que valorizam a escrita dentro de seus padrões, as características do trabalho acadêmico e, principalmente, a construção do conhecimento.

\chapter{Formato} \label{ch:formato}

\section{Digitação do Documento}

\par Os trabalhos devem ser apresentados de modo legível, em folha branca de tamanho A4. O documento deve ser digitado em espaço 1,5 (exceto capa, folha de rosto, resumo e referências bibliográficas/bibliografia, que devem ter espaço um), ocupando apenas o anverso da página. Recomenda-se a utilização da fonte 12 (exceto para capa, folha de rosto e títulos de capítulos, em que deve ser utilizado o tamanho 14 da fonte) em \textit{Times New Roman}, desde que uniforme em todo o documento. Tipos itálicos são usados para nomes científicos e expressões estrangeiras, títulos de obras ou periódicos e citações. A margem à direita deve ser justificada, bem como a da esquerda, exceto no início de parágrafo, cujo recuo deve ser obrigatório de 1,25 cm.

\section{Margens do Documento}

\par As margens devem permitir encadernação e reprodução corretas.
\begin{itemize}
\item Margem esquerda: 3.0 cm
\item Margem direita : 2.0 cm
\item Margem superior : 3.0 cm
\item Margem inferior : 2.0 cm
\end{itemize}

\section{Numeração das Páginas do Documentos}
\par As páginas são contadas \textbf{a partir da folha de rosto} e não são numeradas. A numeração destas acontece sequencialmente \textbf{a partir da Introdução}, em algarismos arábicos, no canto superior direito, sem traços, pontos ou parênteses. As páginas do anexo e/ou apêndice (quando houver) devem ser numeradas de maneira contínua de forma a dar seguimento à paginação do texto principal.

\section{Título e Subtítulo} \label{sec:tituloSubtitulo}

\par É destacado gradativamente, usando-se os recursos de negrito. O título de capítulo deve ser digitado em letra maiúscula, negrito, fonte tamanho 14, sem parágrafo, utilizando-se algarismos arábicos. Os subtítulos devem ser impressos em letra 12 negrito e só a primeira letra das palavras principais em maiúscula.  Todos os capítulos devem ser iniciados em páginas próprias, ainda que haja espaço útil na folha. Este guia pode servir como exemplo. Entretanto, os subtítulos (como a \autoref{sec:tituloSubtitulo} destas orientações) ocupam a mesma página do capítulo, quando há espaço.  Entre o título e o texto, ou entre um subtítulo e o texto, deve-se deixar um \verb*#<#enter\verb*#>#. E entre um subtítulo e outro, após terminar o texto, são dois \verb*#<#enters\verb*#>#.


\chapter{Estrutura} \label{ch:estrutura}

\section{Hipótese e Problematização} \label{sec:hipoteseProblematizacao}

\par A pesquisa científica deve seguir o método científico para receber o qualificativo termo “científico”. A metodologia cientifíca pode ser compreendida de acordo com a \autoref{fig:metodoFluxo}: 

\begin{figure}[!htb]
\caption{Fluxo do Método Científico}
\label{fig:metodoFluxo}
\begin{center}
\includegraphics[height=5cm]{img/JPG/fluxometodo.jpg}\\[0pt]
\end{center}
\legend{Fonte: os autores}
\end{figure}

\par A fase da identificação do problema e do enunciado das hipóteses constitui uma das etapas mais importantes do trabalho visto que direciona e norteia o desenvolvimento da pesquisa científica.

\par Exemplo: Quando você identifica que o telefone da sua casa não está funcionando, inconscientemente, você utiliza-se do método para resolver esse problema, conforme apresentado na \autoref{fig:exemploMetodo}:

\begin{figure}[!htb]
\caption{Exemplo de Aplicação do Método Científico}
\label{fig:exemploMetodo}
\begin{center}
\includegraphics[height=3.5cm]{img/JPG/exemploMetodo.jpg}\\[0pt]
\end{center}
\legend{Fonte: os autores}
\end{figure}

\section{Etapas} \label{sec:etapas}

\par A estrutura do trabalho científico/acadêmico deve apresentar as seguintes etapas mostradas na \autoref{tab:elementos}:

\begin{table}[!htb]
\caption{Elementos do Trabalho Científico}
\label{tab:elementos}
\begin{center}
\begin{tabular}{ |l|l| }
\hline
Estrutura & Elemento \\ \hline
\multirow{5}{*}{Pré-Textuais} & -capa \\
 & -agradecimentos \\
 & -resumo \\
 & -listas(ilustrações, quadros) - quando houver \\
 & -sumário \\ \hline
\multirow{3}{*}{Textuais} & -introdução \\
 & -desenvolvimento \\
 & -conclusão \\ \hline
\multirow{2}{*}{Pós-Textuais} & -referências bibliográficas \\
 & apêndices e/ou anexos (quando houver) \\
\hline
\end{tabular}
\end{center}
\end{table}
 
\par Um trabalho científico é composto de elementos que o caracterizam como tal. São divididos, como se observa na tabela acima, em elementos pré-textuais, textuais, e pós-textuais. Há elementos que são obrigatórios, ou seja, aqueles que propiciam ao leitor as informações importantes sobre o trabalho escrito, e elementos que são opcionais, ou seja, aqueles que trazem informações secundárias ao leitor, permitindo que este obtenha complementos de leitura.

\chapter{Elementos Pré-Textuais} \label{ch:elementos}

\par Os elementos pré-textuais são aqueles que antecedem o texto propriamente dito. Isso significa que, através dos pré-textuais é possível reconhecer o autor do trabalho, a natureza, a localização de seu conteúdo, entre outros aspectos. 

\section{Capa}
\par Deve constar, autoria, título do trabalho, local e data, dispostos em fonte \textit{Times New Roman} 14 (apenas o autor/autores do trabalho deve ser escrito em letra 12) em negrito, espaçamento simples, segundo modelo da página seguinte. No exemplo (\autoref{an:Diagramacao}), observam-se os elementos identificadores do trabalho. O leitor deve reconhecer, pela capa, que o trabalho pertence à determinada instituição, o curso específico, o título, os autores estes em fonte 12 de forma a indicar que contribuíram modestamente com sua pesquisa para o mercado profissional -, o local e a data de realização. Ao elaborar a capa, colocar seu próprio nome no lugar de ``Jorge...'' do exemplo. Caso haja mais de um autor, os nomes podem aparecer em ordem alfabética, elencados um abaixo do outro.

\section{Agradecimentos}

\par Em uma página, no centro colocar o título ``AGRADECIMENTOS" em letra 14, negrito, todas as letras em maiúsculo, com espaçamento duplo. Abaixo, após um \verb*#<#enter\verb*#>#, com o documento justificado e recuo de parágrafo padrão, iniciam-se os agradecimentos propriamente ditos. Obedecem a uma regra de praxe dos documentos científicos, ou seja, presta-se gratidão às pessoas /empresas/instituições envolvidas direta e/ou indiretamente na realização da pesquisa. Isto inclui familiares, orientador, professores, funcionários, chefes e outras pessoas relacionadas aos ambientes profissionais dos autores do trabalho, além dos amigos que contribuíram com o apoio, leitura ou indicação. Citam-se os nomes.

\par Exemplo de agradecimento \footnote{Neste exemplo, mantém-se a redação dos autores.}

\begin{quote}
\begin{center}
\textbf{AGRADECIMENTOS}
\end{center}
\par Agradecemos a todos que tornaram possível a realização deste projeto, especialmente ao ex-presidente da Instituição Grupo Sol da Cidadania Ari Ribeiro e a todos os voluntários, por nos abrir as portas da instituição e nos ceder de boa vontade as informações necessárias para a realização deste projeto.
\par Agradecemos também a todos aqueles da Instituição Prof. Luiz Rosa que nos ajudaram no que foi preciso, alunos, funcionários e professores, especialmente aos Profs. Micaela Deyust dos Santos Pincinato, José Antonio da Cunha Carnevalle e Daiane Pinelli Pavan, que nos orientaram durante toda a construção do projeto. 

\end{quote}

\section{Resumo Em Português}
\par O resumo científico, escrito em língua materna, é um tipo de texto cuja finalidade é sintetizar, ou seja, mostrar objetivamente de que se trata seu trabalho. Descrevem-se, então, os principais tópicos do texto original e indica-se, sucintamente, seu conteúdo. Redigido pelo próprio autor do trabalho, o resumo deve ser a síntese dos pontos relevantes do texto, em linguagem clara, concisa e direta. A palavra ``RESUMO" deve ser escrita em fonte 14, negrito, todas as letras em maiúsculo e alinhamento no centro. No corpo do texto, separado após um \verb*#<#enter\verb*#># deve-se ressaltar o tema, a importância e o objetivo do trabalho. Descreve-se brevemente a realização da pesquisa feita pelo grupo e sua contribuição para com o mercado profissional. O resumo, escrito em 3ª. pessoa do singular e /ou na forma impessoal, o que propicia maior credibilidade ao texto, deve ser redigido em um só parágrafo, sem o recuo de margem, com espaçamento simples, fonte 12. Contém entre 250 e 500 palavras, salientando-se a importância da colocação de três palavras-chave. Estas são aquelas que mais se repetem ou são de grande importância no trabalho como um todo. 

\par Exemplo de Resumo \footnote{Foi mantida a redação original}

\begin{quote}
\begin{center}
\textbf{RESUMO}
\end{center}
\par O presente trabalho estuda o funcionamento de uma biblioteca. O principal objetivo é analisar a infra-estrutura, os equipamentos e os serviços da Biblioteca Leme do Prado. A pesquisa desenvolveu-se a partir de uma revisão de literatura na área de biblioteconomia sobre a temática Biblioteca Universitária. Os dados foram levantados através de formulário implementado e atualizado a partir da obra \textit{Guia das Bibliotecas Brasileiras} publicado pelo IBGE, em 1980, e aplicado à diretora da Biblioteca, permitindo traçar o seu perfil com relação à sua localização, o tamanho das coleções (total de volumes), a existência e adequação da área destinada à Biblioteca, a existência de equipamentos, a quantidade e qualificação de pessoal. Abordou-se o histórico da Biblioteca no Brasil e suas funções. Analisou-se a importância da Biblioteca Universitária na disseminação da informação e na construção de uma sociedade mais democrática e dinâmica.

\par Palavras-chave: Biblioteca, Funções, Bibliografia

\end{quote}

\section{Sumário}

\par Consiste na enumeração dos capítulos do trabalho, na ordem em que aparecem no texto, com a página inicial de cada capítulo. Tem como finalidade facilitar a busca dos assuntos mencionados nos capítulos.
\par A palavra ``SUMÁRIO", centralizada, deve ser escrita em fonte 14, negrito. Deixa-se um espaço em branco de espaçamento 1,5. Os capítulos, escritos em letra maiúscula, devem ser numerados em algarismos arábicos, a partir da Introdução até os Anexos/Apêndices. Havendo subdivisão nos capítulos, deve ser adotada a numeração progressiva, sempre em número arábico. Apenas a primeira letra das palavras do subtítulo escreve-se em letra maiúscula. A \autoref{fig:sumario} apresenta um modelo de sumário.

\begin{figure}[!htb]
\begin{center}
\includegraphics[height=10cm]{img/JPG/sumario.jpg}\\[0pt]
\end{center}
\caption{Exemplo de Sumário. Extraído de \citeonline{Pupo2003}}
\label{fig:sumario}
\end{figure}

\subsection{Listas}
\par ``Listas" (ilustrações ou tabelas / gráficos/ quadros/imagens/ figuras) é um elemento que se destina a identificar os elementos gráficos na ordem em que aparecem no texto, indicando seu título, o número da página em que estão impressos. 

\chapter{Elementos Textuais} \label{ch:pesquisa}
\par Como todos os trabalhos científicos, a organização do texto deve obedecer a uma sequência, ou seja, Introdução, Desenvolvimento e Considerações Finais, dividindo-se em capítulos conforme a natureza do assunto. A parte textual propriamente dita, no Ensino Médio  e Técnico, deve ter no mínimo 10 e no máximo 20 páginas, neste nível de proposta.
\par É importante salientar que um texto acadêmico deve caracterizar-se pela originalidade e autoridade e preservar-se os direitos autorais, sendo proibida a apropriação indevida das obras.

\par Para corroborar nossa ideia, citamos \citeonline{Muller2003}:

\begin{citacao}
\par ``(...) Citar não é pecado nem errado; porém, o que não é permitido é a transcrição de partes do texto, bem como de ideias e sugestões de outro autor, sem se fazerem as referências devidas. "
\par``(...)É obrigatório, pois, ao autor do trabalho, indicar os dados completos das fontes de onde foram extraídas as citações (...)."
\end{citacao}

\par O que foi acima citado esclarece a questão da preservação dos direitos autorais \footnote{A Lei no. 9610/98 garante os direitos autorais e pode ser consultada no site da Associação Brasileira de Direitos Bibliográficos (ABDB) no site \urlstyle{sf}\url{www.abdr.org.br}.} . Tais observações também servem para os textos extraídos da internet: não é permitida, ao aluno-pesquisador, a apropriação indevida destes documentos, bem como de imagens/fotografias sem a menção à fonte.

\par Nas seções seguintes, explicita-se a aplicação de cada elemento textual.

\section{Introdução}
\par Nesta primeira parte do texto os autores devem fazer uma apresentação geral do assunto do trabalho e definir sucinta e objetivamente o tema abordado. Também se acrescenta a justificativa sobre a escolha do tema e métodos empregados. Todo trabalho de pesquisa segue uma metodologia, que abrange desde visitas, passando por questionários, entre tanto outros métodos. Deve-se também colocar os objetivos de pesquisa, com especificação sobre os aspectos que serão ou não abordados. Enfim, é feito um resumo de cada capítulo do trabalho, em diversos parágrafos. A introdução deste livreto segue o padrão acadêmico e pode ser tomada como exemplo e sugestão.
\par Abaixo, encontram-se exemplos, que devem ser adaptados de acordo com as ideias referentes ao tema do seu trabalho. Importante salientar que em itálico indica apenas o significado do conteúdo do parágrafo que serve de exemplo. O que está escrito em itálico não deve fazer parte da introdução do seu trabalho, já que esta deve compor-se como um texto corrido, dividido em parágrafos

\section{Desenvolvimento}

\par É o trabalho propriamente dito, ou seja, é o texto que contém o detalhamento da pesquisa, citações para explicar um determinado pensamento do grupo, explicação da importância do assunto estudado e os valores agregados. Pode conter também uma análise dos prós e contras para convencimento do leitor.
\par Para organização da escrita acadêmica, é fundamental a divisão do desenvolvimento em capítulos. Cada um deles pode descrever de onde partiu a pesquisa, informar o nome e dados da empresa pesquisada, a fundamentação teórica, a coleta de dados e análise dos resultados obtidos, de modo que o texto obedeça a uma sequência. A fundamentação teórica é a exposição da teoria defendida com base em argumentos de autoridade, ou seja, de autores renomados que possibilitam, através de seus estudos, entender o universo de sua área de pesquisa. Já a coleta de dados refere-se à busca de informações sobre o objeto de estudo, as quais serão organizadas de acordo com a intenção do autor ou autores do trabalho. A partir disso, as informações serão analisadas tendo como orientação a teoria defendida anteriormente. Assim estruturados, os capítulos devem propiciar maior compreensão, por parte do leitor, das ideias do grupo sobre o Projeto Integrador.
\par Os quadros e as tabelas são conjuntos de informações apresentados em colunas e linhas. Nos quadros as informações expostas são predominantemente palavras e textos. Já as tabelas apresentam informações predominantemente quantitativas, em forma de números, valores, sejam ou não tratados estatisticamente. A indicação do quadro/tabela deve ocorrer diretamente no texto, fazendo parte de uma afirmação isolada ou entre parênteses.  

\par Os quadros e as tabelas devem estar próximos do trecho ao qual fazem referência. Eles devem estar centralizados na página e com disposição de texto alinhada. Caso uma tabela, por exemplo, não caiba em uma página, escrever embaixo a palavra (em maiúsculo) ``CONTINUA" (em letra 10) e continuar a tabela na outra página (repita o título da tabela, seguido da palavra ``CONTINUAÇÃO" - também em maiúsculo - em letra 10).
\par Ilustrações são todas as figuras, desenhos, rascunhos, esquemas, diagramas, fluxogramas, organogramas, fotografias, mapas e demais imagens da dissertação, ou tese, que servem para elucidar, explicar, complementar e simplificar o entendimento do texto. Sua indicação deve ocorrer diretamente no texto, fazendo parte de uma afirmação isolada ou entre parênteses, assim como ocorre com os quadros e tabelas. 

\par Os gráficos e demais ilustrações devem estar próximos do trecho ao qual fazem referência. Eles devem estar centralizados na página e com disposição de texto alinhada. Todas as ilustrações do desenvolvimento da obra devem se enquadrar nas margens adotadas, não podendo ultrapassá-las em nenhuma hipótese. 
\par Os textos e números inseridos nos gráficos e nas ilustrações devem ser redigidos em Times New Roman tamanho 12 na cor preta e com espaçamento simples. Cores são facultativas, mas sua utilização deve ser moderada. Os demais detalhes estéticos ficam a cargo do autor, desde que a configuração seja mantida em todo o trabalho. Uma legenda correspondente ao gráfico ou demais ilustrações deve ser obrigatoriamente indicada. A legenda (tamanho 12) de ambos os casos deve apresentar o número sequencial dentro da obra, o título e a fonte (letra desta tamanho 10) abaixo do objeto.

\section{Considerações Finais}
\par É o espaço de retomada do tema e dos assuntos e resultados da pesquisa. Isso significa que uma síntese é elaborada pelo grupo, seguida de questionamentos e propostas para projetos futuros. Também deve caracterizar-se pela escrita própria, na qual se evita a cópia de textos alheios. Também é nas considerações que se comparam os objetivos iniciais e os finais. Isto é, no início do trabalho apresenta-se uma meta que vai se confirmar ou não conforme o andamento da pesquisa, ao final do processo.

\chapter{Mecanismos de Fidelidade do Texto: Citações e Notas de Rodapé}

\section{Apresentação de Informação no Texto - Citação}

\par Acreditamos ser importante destacar um comentário sobre a citação e suas implicações: as citações quando relacionadas diretamente ao texto, têm o poder de enriquecer a matéria, pois denota pesquisa e preocupação por parte do autor do trabalho com relação aos autores escolhidos, geralmente relevantes para o assunto \cite{Muller2003}.
\par Entende-se citação como a menção no texto de uma informação colhida de outra fonte, cujo objetivo é fundamentar as ideias defendidas no trabalho acadêmico. Salienta-se que uma citação solta num trabalho acadêmico, sem a devida relação com aquilo que se está escrevendo, denota pobreza de escrita. Tal relação no texto pode ser elaborada tendo em vista comparações, contra-argumentações, exemplificações, enumerações, definições, interrogações e informações históricas, entre outras.  

\par Há que se identificar o autor da citação para que o trabalho não se caracterize como cópia pura. Essa identificação segue normas rigorosas padronizadas pela ABNT e adaptadas por este manual.

\begin{alineas} %for small alpha-characters within brackets.
\item Autor
\par Indicação do Sobrenome do autor com inicial em maiúscula, seguido da data.
\par Exemplo: \citeonline{ROSA2006} enfatiza que a interpretação do texto.
\item Dois ou Três Autores
\par Indicação dos Sobrenome do autor com inicial em maiúscula seguido da data.
\par Exemplo: \citeonline{Santos2008} definem a pesquisa científica como... 
\item Mais de Três Autores
\par Indicação do primeiro autor, seguido da expressão \textit{et al.} acrescido da data.
\par Exemplo: \citeonline{Medeiros2003} remete-se à questão da leitura e ao início do mundo da escrita...
\end{alineas}

\par A lista final de referências bibliográficas e de bibliografia deve ter arranjo alfabético.

\par As citações podem ser divididas em citações diretas, indiretas e citação da citação. Salienta-se que para a elaboração dos Projetos Integradores serão utilizados os tipos: citação direta e citação indireta.

\subsection{Citação Direta}
\par É a cópia exata ou transcrição literal de outro texto. O tamanho de uma citação determina sua localização no texto e pode ser utilizada da seguinte forma:

\begin{alineas}
\item Até três linhas deve ser incorporada ao parágrafo, entre aspas duplas e em itálico, de duas formas. 

\par Exemplos:

\par A leitura está presente em nossa vida de tal forma que sempre nos lembramos dos livros que lemos quando pequenos.``\textit{ Talvez não haja na nossa infância dias que tenhamos vivido tão plenamente como aqueles que (...) passamos na companhia de um livro preferido.}"\cite{Pupo2003}.

Todos esses elementos, dependendo da forma como são conduzidos, implicarão positiva ou negativamente na formação do leitor proficiente. Para \citeonline{Lima2008}: ``\textit{são os elementos relevantes ou representativos os que contam, em função do significado do texto, a experiência do leitor é indispensável para construir o sentido (...)}". 

\item citação mais longa, acima de 4 linhas, inclusive, deve figurar abaixo do texto, alinhado à direita, com recuo de 4 cm a partir da margem inicial esquerda de 3 cm, espaçamento simples,  letra 10 (sem itálico, sem aspas, sem negrito). 
\par Exemplo:
\par Dessa maneira, deve-se considerar na construção de sentido a existência de um texto para que o leitor possa preencher-lhe os vazios existentes, dada a relação móvel que o texto tem com o leitor na dependência do reconhecimento do seu gênero como sustenta Chartier: 

\begin{citacao}

Os textos não são depositados nos objetos, manuscritos ou impressos, que o suportam como receptáculos, e não se inscrevem no leitor como o fariam em cera mole. Considerar a leitura como um ato concreto requer que qualquer processo de construção de sentido, logo de interpretação, seja encarado como estando situado no cruzamento entre [...] leitores dotados de competências específicas [...] e [...] textos cujo significado se encontra dependente dos dispositivos discursivos e formais ... \cite{Chartier1999}
\end{citacao}

\end{alineas}

\subsection{Citação Indireta}

\par É a expressão da ideia contida na fonte citada, sem transcrição, dispensando o uso de aspas duplas. Menciona-se, de qualquer forma, o nome do autor parafraseado e a data de publicação da obra. 

\par Exemplo:

\begin{quote}
\par Para \citeonline{MAZZETTI2001}, a leitura é questão de suma importância tanto que sua falta tende a ser classificada como uma aflição. Sua preocupação, porém, refere-se a: a) o fato dos elementos que afetam a leitura e b) os danos que ela provoca, nomeando esses fatores como maldição ou praga. (...)
\end{quote}

\par Faz-se necessário colocar uma observação acerca das citações de textos online. Quando é um artigo, por exemplo, com autoria, faz-se a seguinte menção: De acordo com Moura (2001, online), a citação é um modo de..... Em caso de textos sem autoria, cita-se o site e, em nota de rodapé, o ideal é apontar o caminho completo do site em questão. Por exemplo: Segundo o site da UNESCO (e em nota de rodapé, indica-se o caminho do site: \verb*#<#\urlstyle{sf}\url{http://www.brasilia.unesco.org/noticias/ultimas/unesco-lanca-curso-para-capacitacao-de-profissionais-de-museus}\verb*#># ), há oferecimento de 160 vagas para o curso de capacitação........

\subsection{Notas de Rodapé}

\par As notas de rodapé \protect\footnote{Algumas informações deste tipo de mecanismo foram adaptadas do site \textit{\protect \urlstyle{sf}\url{http://www.admbrasil.com.br/abnt.htm.}}} destinam-se a prestar esclarecimentos ou tecer considerações que não devam ser incluídas no texto para não interromper a sequência lógica da leitura. 

\par Para se fazer a chamada das notas de rodapé, usam-se algarismos arábicos, na entrelinha superior sem parênteses, com numeração consecutiva para cada capítulo ou parte, evitando-se recomeçar a numeração a cada página. Hoje, o documento do Word conta com a janela ``Inserir". Clicando-se nela abre-se em uma outra janela ``Referências", ``Nota de Rodapé", e a chamada é feita automaticamente.


% ----------------------------------------------------------
% Finaliza a parte no bookmark do PDF
% para que se inicie o bookmark na raiz
% e adiciona espaço de parte no Sumário
% ----------------------------------------------------------
\phantompart

\chapter{Elementos Pós-Textuais} \label{ch:posTextuais}

\par São aqueles que contribuem na concepção e realização do trabalho acadêmico e que necessitam ser informados, segundo uma apresentação própria definida pela ABNT, sem, entretanto, integrarem o corpo do texto propriamente dito. Também representam o conjunto de elementos que permitem a identificação, no todo ou em parte, de documentos impressos ou registrados em diversos tipos de materiais utilizados para a ampliação do conhecimento do autor do trabalho acadêmico. Vale acrescentar que deve haver um equilíbrio entre as fontes pesquisadas e referenciadas, bem como um maior número de livros/revistas científicas do que textos online. Selecione as fontes online mais relevantes e confiáveis e habitue-se a fundamentar boa parte de sua pesquisa em livros (a não ser que seu trabalho seja muito específico e relacionado/restrito à leitura de textos online).


\section{Referências Bibliográficas}

\par A referência bibliográfica caracteriza-se como textos de apoio citados de forma direta e/ou indireta no trabalho acadêmico e que garantem a credibilidade da produção textual. Isso significa que o aluno-pesquisador, ao fazer citações de outras fontes em seu documento e ao indicar a autoria delas, demonstra ao leitor honestidade científica. 

\par Assim, ao final do seu trabalho, o aluno-pesquisador deve listar as referências utilizadas no corpo do texto. Essa identificação segue um padrão convencionado pela ABNT, adaptado para a apresentação dos trabalhos acadêmicos e para as monografias dos Programas de Pós-Graduação.

\subsection{Livros}

\par Os elementos a serem considerados são: autor ou autores; título em itálico; subtítulo, se houver; edição, partindo-se da 2ª; local de publicação; editora; ano de publicação; 
\par Se alguma informação não constar da obra pesquisada, podem-se utilizar alguns recursos que auxiliam na composição dos elementos. São eles:

\begin{alineas}
\item se o título ou subtítulo da obra for muito longo, utilizam-se reticências após a escrita das primeiras palavras.
\item acréscimos à edição como \textit{revista / ampliada} devem ser mencionados de forma abreviada.
\item se o local da publicação não estiver mencionado na obra, é válida a informação [S.l.], \textit{Sine loco}.
\item a obra pesquisada que não mencionar a editora em suas partes pode ser referenciada na sua ordem pela denominação [s.n.],\textit{sin nomine}.
\item não havendo informações referentes ao local de publicação nem à editora, recorre-se à denominação [S.l.:s.n.].
\item quando o ano de publicação não for mencionado, pode-se informar uma data aproximada, com base em alguma orientação mencionada na obra, como exemplo: [1981].
\item a indicação das páginas ao final da referência pode ser feita de forma simples: p. 17-25 ou p. 31-6 (da página 31 à página 36). Quando não houver paginação ou esta for irregular, as menções se apresentam como: paginação irregular ou não paginado.
\end{alineas}

\par Exemplos:

\subsubsection{Com um autor:}
\begin{quote}
\par KLEIMAN, Angela. \textit{Texto e leitor: aspectos cognitivos da leitura.} 7ª. ed. Campinas, SP: Pontes, 2000. 82p.
\end{quote}

\subsubsection{Com dois autores:}
\begin{quote}
\par FIORIN, José Luiz; SAVIOLI, Francisco Platão. \textit{Para entender o texto: leitura e redação}. 16ª. ed. São Paulo: Ática: 2003. 431p.
\par \verb*#______________________________________#. \textit{Lições de texto: leitura e redação.} São Paulo: Ática, 1996. 432p. 
\par O tracejado utilizado na bibliografia ou nas referências bibliográficas indica que a obra é do autor listado anteriormente.
\end{quote}


\subsubsection{Com três autores:}
\begin{quote}
\par SCHERMERHORN (Jr.), John R.; HUNT, James G.; OSBORN, Richard N. \textit{Fundamentos de Comportamento Organizacional.} 2a. ed. Porto Alegre: Bookman, 1999. 328p.
\end{quote}

\subsubsection{Com quatro ou mais autores:}
\begin{quote}
\par PAULINO, Graça \textit{et al}. \textit{Tipos de Textos, Modelos de Leitura.} Belo Horizonte: Formato, 2001. 163p.
\end{quote}


\par Há que se mencionar aqui algumas observações: 
\begin{alineas}
\item com quatro autores ou mais destacam-se todos os nomes quando isto for indispensável. 
\item nas obras em que reúnem vários trabalhos de autores diferenciados, a entrada se dá pelo organizador ou coordenador do trabalho. 

\par Exemplo:

\par ABREU, Márcia (Org.). \textit{Leitura, história e história da leitura.} Campinas, SP: Mercado das Letras; Associação de Leitura do Brasil; São Paulo: Fapesp, 1999. 640p. (Coleção Histórias de Leitura) 

\item Quando não se conhece o autor, a referência inicia-se com o título da obra.
\end{alineas}


\subsection{Capítulos de Livros:}

\par Autor ou autores da parte em referência, título da parte e subtítulo, (se houver), expressão In: autor ou autores da obra no todo quando estes forem mencionados ou uso de travessão (na proporção de seis caracteres) caso não se apresentem os autores,  título da obra no todo (em itálico) e subtítulo (se houver), edição, editora, ano da publicação e localização.

\par Exemplo: 

\begin{quote}
CHARTIER, Roger. As Revoluções da Leitura no Ocidente. In: ABREU, Márcia (Org.). \textit{Leitura, história e história da leitura.} Campinas, SP: Mercado das Letras; Associação de Leitura do Brasil; São Paulo: Fapesp, 1999. p. 19-31. (Coleção Histórias de Leitura) 
\end{quote}

\subsection{Trabalhos Monográficos}

\par Podem ser mencionados no todo ou em partes (capítulo ou seção). 
\par A sequência apresenta-se dessa forma: autor, título (itálico), subtítulo (se houver),  ano,  número total das páginas, o tipo de trabalho monográfico, nível e área do curso, nome da instituição e local.

\par Exemplo:

\begin{quote}
VENTURA PUPO,  Lucí Bertolo. \textit{Práticas de Leitura em Alunos do Ensino Médio.} 2003. 104 p. Dissertação (Mestrado em Educação) - Universidade São Francisco, Bragança Paulista.
\end{quote}


\subsection{Obra de Entidade}
\begin{quote}
\par UNIVERSIDADE SÃO FRANCISCO. \textit{Catálogo de Publicações 2001}. Bragança Paulista, SP, 2001.

\par Se a entidade em questão apresentar nome mais genérico, dá-se a entrada por meio do órgão superior ou pela unidade geográfica da qual faz parte.

\par Exemplo:

\par BRASIL. Ministério de Minas e Energia. Departamento de Administração. \textit{Relatório de Atividades}, 1990-1995. Brasília, 1996.

\end{quote}

\par A referência de parte do texto monográfico apresenta a ordem: autor ou autores do texto em referência, o título da parte lida, expressão In:, travessão (na proporção de seis caracteres), título da obra no todo (itálico), subtítulo (se houver), ano, localização do texto referido, o tipo de trabalho monográfico, nível e área do curso, nome da instituição e local.

\par Exemplo:

\begin{quote}
VENTURA PUPO, Lucí Bertolo. Considerações dentro da Complexidade da Leitura. In: \textit{Práticas de Leitura em Alunos do Ensino Médio.} 2003. p. 6-40. Dissertação (Mestrado em Educação) \verb*#-# Universidade São Francisco, Bragança Paulista.
\end{quote}


\subsection{Outras formas de referência}

\subsubsection{Acesso \textit{On-Line}}

\par Os elementos informativos para livro devem seguir a sequência: autor ou autores (se houver), título (itálico), subtítulo (se houver), ano da publicação (se houver), expressão Disponível em:, endereço eletrônico entre < >, expressão Acesso em:, data do acesso.

\par Exemplo de acesso \textit{On-Line} \protect \footnote{Mais indicações sobre como citar textos obtidos em meio eletrônico no site <\protect \urlstyle{sf}\url{http://www.quatrocantos.com/tec_web/refere/INDEX.HTM}> Acesso em 19/01/2009}:

\begin{quote}
JOÃO XXIII. \textit{Mater et magistra: evolução da questão social à luz da doutrina cristã.} Roma. 1961. Disponível  em: \protect \urlstyle{sf}\url{http://www.vatican.va/holy_father/john_xxiii/encyclicals/documents/hf_jxxiii_enc_15051961_mater_po.html}. Acesso em: 20 nov. 2002.
\end{quote}

\par As informações de referência para periódicos em meio eletrônico devem seguir a sequência: autor ou autores do texto quando houver, o título do texto, subtítulo (se houver), o título do periódico (em itálico), local da publicação, data, seção ou parte do periódico, páginas, descrição do objeto (CD ou \textit{pen drive}) ou expressão "Disponível em:"; endereço eletrônico entre < >; expressão "Acesso em:"; data do acesso. Os trabalhos monográficos apresentam a ordem: autor; título (itálico); subtítulo,se houver; ano; número total das páginas; o tipo de trabalho monográfico; nível e área do curso; nome da instituição; local; expressão Disponível em:; endereço eletrônico entre < >; expressão Acesso em:; data do acesso.

\subsubsection{Publicação Periódica (Revista, Boletim)}

\par Deve-se mencionar o autor do texto quando houver; o título do texto; o título do periódico (em itálico); local da publicação; número do volume ou o ano; número do fascículo; páginas em que se encontra o texto; data do fascículo.
\par Exemplo:
\begin{quote}
Como trabalhar por competências. \textit{Nova Escola} - Especial Ensino Médio, São Paulo, ago. 2002, Caderno Especial.
\end{quote}

\subsubsection{Publicação Periódica (Artigo ou Matéria de Jornal)}
\par Citam-se os seguintes elementos informativos: o autor ou autores do texto quando houver; o título do texto; subtítulo, se houver; o título do periódico (em itálico); local da publicação; data; seção ou parte do periódico; páginas.
\par Exemplo:
\begin{quote}
Provão da desigualdade. \textit{Folha de São Paulo}, São Paulo, 12 dez. 1999. Editorial, p. 2
\end{quote}

\subsubsection{Legislação}

\par É possível referenciar os elementos legislativos conforme exemplo abaixo: 
\begin{quote}
BRASIL. Decreto-lei n. 2423, de 7 de abril de 1988. Estabelece critérios para pagamento de gratificações e vantagens pecuniárias aos titulares de cargos e empregos da Administração Federal direta e autárquica. \textit{Diário Oficial da União}, Brasília, v.126, n.66, p. 6009, 8 abr. 1988. Seção 1, pt.1.
\end{quote}

\section{Anexo e Apêndice}
\par O primeiro é um elemento não elaborado pelo autor, que documenta, esclarece, prova ou confirma as ideias expressas na parte textual do trabalho acadêmico, por exemplo: modelos de documentos utilizados em empresas, leis e pareceres que complementam uma idéia expressa no desenvolvimento do trabalho, dentre outros.

\par O apêndice é um elemento que consiste em um documento ou texto elaborado pelo autor, com o intuito de complementar sua argumentação sem prejuízo do trabalho. Podem-se citar, como exemplos, o cronograma estabelecido para o cumprimentos de prazos, modelos de questionários, o roteiro de uma entrevista, dentre outros.

% ----------------------------------------------------------
% ELEMENTOS PÓS-TEXTUAIS
% ----------------------------------------------------------
\postextual
% ----------------------------------------------------------

% ----------------------------------------------------------
% Referências bibliográficas
% ----------------------------------------------------------
\bibliography{bib/referencias}

% ----------------------------------------------------------
% Anexos
% ----------------------------------------------------------

% ---
% Inicia os anexos
% ---
\begin{anexosenv}

% Imprime uma página indicando o início dos anexos
\partanexos

% ---
\chapter{Diagramação de Página do Trabalho Científico} \label{an:Diagramacao}

\begin{figure}[!htb]
\caption{Diagramação das páginas do trabalho científico.}
\label{fig:anexo}
\begin{center}
\includegraphics[height=15cm]{img/JPG/anexoa.jpg}\\[0pt]
\end{center}
\end{figure}

\begin{figure}[!htb]
\caption{Folha de Rosto}
\label{fig:anexoB}
\begin{center}
\includegraphics[height=15cm]{img/JPG/anexob.jpg}\\[0pt]
\end{center}
\end{figure}

\end{anexosenv}

%---------------------------------------------------------------------
% INDICE REMISSIVO
%---------------------------------------------------------------------
\phantompart
\printindex
%---------------------------------------------------------------------

\end{document}
